Phase I is a discovery study: its role is to establish the phenomenon, the
taxonomy, and the existence of task-side headroom, and to motivate the
confirmatory design. None of its results is confirmatory: 24 of its 151
evaluation items overlapped protocol-development measurements, the human
control was designed after inspecting development traces, and its judge was
a conservative proxy (Appendix~\ref{app:phase1-provenance}). We report it
with those limits stated.

\subsection{Harnesses, outcomes, and the three metrics}\label{sec:metrics}
A \textbf{harness} $H$ wraps a frozen target model $M$: it receives the
natural-language task $x$, issues one or more calls to $M$, optionally
executes intermediate artifacts (e.g.\ SQL), and returns a final answer
scored programmatically. With $H_0$ the bare single-call baseline and
$Y \in \{0,1\}$ correctness, the \textbf{harness marginal value} is
$\Delta(x,H,M) = Y(x,H,M) - Y(x,H_0,M) \in \{-1,0,+1\}$
($+1$ repair, $-1$ harness-induced harm).

We propose that harness-generation and
harness-routing pipelines report three precondition metrics on a held-out
set, all computable from the outcome matrix any evaluation already produces:
\textbf{(i) pairwise outcome disagreement} --- over all harness pairs, the
fraction of tasks where correctness differs; \textbf{(ii) union repair
rate} --- the fraction of bare errors some population member repairs;
\textbf{(iii) bare-inclusive oracle headroom} --- over held-out tasks
$x_1,\dots,x_n$,
\[
H(P) \;=\; \frac{1}{n}\sum_{i=1}^{n} \max_{H \in P \cup \{H_0\}} Y(x_i, H)
\;-\; \max_{H \in P \cup \{H_0\}} \frac{1}{n}\sum_{i=1}^{n} Y(x_i, H).
\]
Positive observed headroom records item-dependent successes in that execution
matrix; independent repeats are needed to establish a stable routing signal.
Disagreement without repair adds no observed coverage beyond bare; repair
without headroom adds no observed oracle advantage over the best fixed member.

\subsection{The collapse, quantified}\label{sec:collapse}
The domain is BIRD text-to-SQL (SQLite; execution-accuracy judging), target
frozen at \textbf{GLM-5.3-Flash} (temperature 0) throughout Phase I. On the
day-1 population --- 14 harnesses (\bare, \react, and 12 evolution
candidates) $\times$ 60 tasks --- average pairwise outcome disagreement is
\textbf{3.3\%} and \textbf{36 of 91 pairs are identical on every task};
among the 12 generated candidates alone (66 candidate pairs; mean pairwise
code similarity 0.143 under the normalized-character metric of the original
code-difference check), disagreement is 3.8\% with 21 identical pairs.
Every candidate had passed a code-difference check (``all files distinct,
non-trivial diffs''). Union repair over bare errors is 7.7\% and oracle
headroom 3.3\,pp; \react{} differs from \bare{} on zero tasks.

Figure~\ref{fig:phenomenon} (Appendix~\ref{app:phase1-provenance}) makes the central relation quantitative: across
all 91 pairs, source-code similarity under the figure's token-sequence
metric has Spearman $\rho = -0.010$ with outcome disagreement --- code
similarity shows \emph{no observed rank association} with behavioral
similarity here. (Pairs overlap; $\rho$ is descriptive, not inferential.)

\subsection{Trace taxonomy: why code diversity fails to become behavioral
diversity}\label{sec:taxonomy}
Auditing execution traces (the sequence of LLM calls and artifact executions
per task) decomposes ``different code'' into three failure classes:
\begin{itemize}
\tightlist
\item \textbf{T1 --- mechanism described, not executed (6/12 candidates).}
The harness issues exactly one LLM call and never executes an artifact ---
the identical execution path as \bare{}. Docstrings claim retrieval,
self-checking, or token caps; none exists in the control flow.
\item \textbf{T2 --- mechanism present but untriggered (\react, 10/10 audited
tasks).} The first SQL executes successfully, so error repair is not
triggered and the final SQL equals \bare's. This does not show that repair
feedback lacks a causal path when a failure occurs.
\item \textbf{T3 --- broken multi-turn path (1/12).} The one candidate
implementing genuine multi-turn repair returns prose as final SQL,
converting bare-correct items into harness-wrong.
\end{itemize}
The taxonomy is actionable: T1 demands an execution-path acceptance gate,
T2 demands controlled faults to test trigger and feedback use, T3 demands
executable-output validation --- all generation-protocol requirements.

\subsection{Human positive control: the headroom exists}\label{sec:positive}
To probe whether distinct hand-written mechanisms can produce additional
observed coverage, we wrote four harnesses by hand with forcibly
distinct execution paths (\texttt{hpc\_repair}: execute-and-repair;
\texttt{hpc\_vote3}: 3-vote on execution results; \texttt{hpc\_schema}:
schema-linking; \texttt{hpc\_hint}: hint-first reordering). Same frozen
target, tasks, and judge as the collapsed population: oracle headroom
\textbf{6.67\,pp} on 60 tasks and \textbf{7.28\,pp} {[}3.3, 10.0{]} on the
151-item set, with repair/schema/hint each producing repairs the others
miss. One disclosed caveat: the control was designed \emph{after}
inspecting day-1 traces, so its headroom is descriptive, not blind. The
bare baseline itself flipped 3 of 60 items between the two day-1 runs
(API nondeterminism at temperature 0).
